\chapter*{Logistic Regression}
\addcontentsline{toc}{chapter}{Logistic Regression}

\section*{A Probabilistic View of Machine Learning}

Up to this point, we have largely viewed classification through a geometric or functional lens (e.g., finding a hyperplane that separates classes). Logistic Regression introduces a \textbf{probabilistic formulation} of learning.

Instead of just predicting a discrete label $y \in \{0, 1\}$ given an input $\xx$, we want to learn the \textbf{conditional probability density} that the label is $y$, given $\xx$. We want our model to output:
$$P(y \mid \xx)$$

\begin{description}
    \item[Generative Models:] These models attempt to learn the joint probability $P(\xx, y)$ or the conditional distribution of the features given the class, $P(\xx \mid y)$. They essentially try to model exactly how the data for each class was generated (e.g., Naive Bayes).
    \item[Discriminative Models:] These models bypass modeling the underlying data distribution and attempt to learn the conditional probability $P(y \mid \xx)$ directly. \textbf{Logistic Regression is a discriminative model.}
\end{description}

\section*{The Logistic Sigmoid Function}

A standard linear model calculates a raw score $z = \mathbf{w}^T\xx + w_0$, which can range from $-\infty$ to $+\infty$. To interpret this output as a probability, we need to map it to the range $[0, 1]$. We achieve this using the Logistic Sigmoid function.

\begin{definition}[Logistic Sigmoid Function]
The sigmoid function, denoted as $\sigma(z)$, "squashes" any real-valued number into the range $(0, 1)$:
$$\sigma(z) = \frac{1}{1 + e^{-z}}$$
\end{definition}

\textbf{Key Properties of $\sigma(z)$:}
\begin{itemize}
    \item As $z \rightarrow \infty$, $\sigma(z) \rightarrow 1$.
    \item As $z \rightarrow -\infty$, $\sigma(z) \rightarrow 0$.
    \item When $z = 0$, $\sigma(0) = 0.5$.
    \item Symmetry property: $\sigma(-z) = 1 - \sigma(z)$.
\end{itemize}

\section*{The Logistic Regression Model}

In binary Logistic Regression (where $y \in \{0, 1\}$), the hypothesis function $f_{\mathbf{w}}(\xx)$ models the probability that the positive class ($y=1$) is true:

$$\hat{y} = f_{\mathbf{w}}(\xx) = P(y=1 \mid \xx) = \sigma(\mathbf{w}^T\xx + w_0) = \frac{1}{1 + e^{-(\mathbf{w}^T\xx + w_0)}}$$

Consequently, the probability of the negative class is simply $P(y=0 \mid \xx) = 1 - P(y=1 \mid \xx)$.

\begin{remark}[The Decision Boundary]
To make a final hard classification, we typically apply a threshold of $0.5$. We predict $y=1$ if $P(y=1 \mid \xx) \geq 0.5$.
Because $\sigma(z) \geq 0.5$ exactly when $z \geq 0$, our prediction rule simplifies to:
$$\text{Predict } 1 \text{ if } \mathbf{w}^T\xx + w_0 \geq 0$$
Therefore, despite the non-linear sigmoid transformation yielding probabilities, \textbf{the underlying decision boundary separating the classes is strictly linear} (a flat hyperplane).
\end{remark}

\section*{Training: Maximum Likelihood and Log-Loss}

Because Logistic Regression outputs probabilities, we do not train it using the standard sum of squared errors. Instead, we use \textbf{Maximum Conditional Likelihood Estimation (MCLE)}. We want to find the parameters $\mathbf{w}$ that maximize the probability of observing the labels in our training dataset.

Instead of maximizing the raw likelihood (which leads to numerical underflow due to multiplying many small probabilities), we typically \textbf{minimize the negative log-likelihood}, which gives us our cost function.

\begin{definition}[Binary Cross-Entropy / Log Loss]
For a dataset of $N$ examples, the Logistic Regression cost function $L(\mathbf{w})$ is defined as:
$$L(\mathbf{w}) = -\frac{1}{N} \sum_{i=1}^N \left[ y^{(i)} \log(\hat{y}^{(i)}) + (1 - y^{(i)}) \log(1 - \hat{y}^{(i)}) \right]$$
where $\hat{y}^{(i)} = \sigma(\mathbf{w}^T\xx^{(i)} + w_0)$ is the predicted probability.
\end{definition}

This loss function heavily penalizes the model if it confidently predicts the wrong class (e.g., predicting $\hat{y} \approx 0$ when the true label $y = 1$ causes the $-\log(\hat{y})$ term to approach infinity).

\section*{Optimization and Regularization}

Unlike Linear Regression, setting the gradient of the Logistic Regression loss function to zero \textbf{does not yield a closed-form analytical solution} due to the non-linear $e^{-z}$ term in the sigmoid.

\begin{description}
    \item[Optimization:] We must use iterative optimization techniques, such as \textbf{Gradient Descent} (or Gradient Ascent, if maximizing likelihood directly). Fortunately, the Log-Loss function is strictly \textbf{convex}, meaning Gradient Descent is guaranteed to find the global optimal weights $\mathbf{w}^*$.
    \item[Regularization:] If the training data is perfectly linearly separable, the algorithm will try to push the probabilities to exactly $1.0$ and $0.0$. To do this, the weights $||\mathbf{w}||$ will grow towards infinity, leading to severe overfitting. To prevent this, we must add $L_1$ (Lasso) or $L_2$ (Ridge) regularization penalties to the loss function.
\end{description}
